\chapter{工程仿真应用案例}

\intro{
	在科学上，没有一件事是和实际应用无关的。
	
	\rightline{——Louis Pasteur}
}

\section{案例一：指数衰减系统的数值仿真}

\subsection{问题描述}

考虑一阶线性微分方程：
\[
\frac{dy}{dt} = -a y, \quad y(0) = y_0, \quad a > 0
\]
该方程刻画了多种物理过程：放射性衰变、RC 电路放电、Newton 冷却定律、以及一级化学反应等。

\subsection{解析解}

分离变量后积分得精确解：
\[
y(t) = y_0 e^{-a t}
\]

\subsection{数值求解——显式 Euler 法}

使用显式 Euler 格式：$y_{n+1} = y_n - a h y_n = (1 - ah) y_n$。

递推 $N$ 步后：$y_N = y_0 (1 - ah)^N$。数值稳定条件为 $|1 - ah| < 1$，即：
\begin{myformula}
	h < \frac{2}{a}
\end{myformula}

当 $h > 2/a$ 时，数值解将发散或振荡，这是显式 Euler 法\textbf{条件稳定}的典型表现。

\begin{example}{对比不同步长}
	取 $a = 1$，$y_0 = 8$，$t_{\max}= 10$。步长 $h$ 从 $0.05$ 到 $0.50$，观察数值解与解析解的偏差。
	
	结论：步长越小、精度越高，但计算量越大；步长超过稳定极限时解发散。
\end{example}

\subsection{MATLAB ode45 实现}

使用 \texttt{ode45} 自适应步长求解器可兼顾精度与效率，自动选择最优步长而无需人工指定。

\section{案例二：等强度柱的 ODE 建模与仿真}

\subsection{物理背景}

考虑一根受拉力的垂直柱（或杆），除顶部的外力 $F$ 外，还需承受自身重量。若采用等截面设计，底部应力最大而顶部应力较小，造成材料浪费。\textbf{等强度设计}的目标是：使每个截面的应力均等于许用应力 $[\sigma]$。

\subsection{数学模型}

截取微元体 $dx$，由静力平衡条件：
\[
\sigma(x) A(x) + dW = \sigma(x+dx) A(x+dx)
\]
其中 $dW = \rho g A(x) dx$ 为微元自重，$\sigma(x) = \sigma(x+dx) = [\sigma]$（等强度要求）。于是：
\[
[\sigma] (A(x+dx) - A(x)) = \rho g A(x) dx
\]
取极限得面积分布的微分方程：
\begin{myformula}
	\frac{dA}{dx} = \frac{\rho g}{[\sigma]} A(x)
\end{myformula}

边界条件：在 $x = 0$（顶部），$A(0) = F / [\sigma]$。解得：
\begin{myformula}
	A(x) = \frac{F}{[\sigma]} \exp\left( \frac{\rho g}{[\sigma]} x \right)
\end{myformula}

\begin{remark}
	底部截面积 $A(L) = A(0) e^{kL}$，其中 $k = \rho g / [\sigma]$。面积沿杆长呈指数增长，这解释了高塔建筑底部必须加粗的原因（如 Eiffel 铁塔）。
\end{remark}

\subsection{总伸长量}

由 Hooke 定律 $\varepsilon = \sigma / E$，总伸长量：
\[
\Delta L = \int_{0}^{L} \varepsilon(x) \, dx = \int_{0}^{L} \frac{[\sigma]}{E} \, dx = \frac{[\sigma] L}{E}
\]
与等截面杆不同，等强度柱的伸长量与截面形状无关，仅取决于许用应力和杆长。

\subsection{数值仿真步骤}

\begin{enumerate}
	\item 参数设置：$F, [\sigma], \rho, E, g, L$；
	\item 计算顶部截面积 $A_0 = F / [\sigma]$ 和指数系数 $k = \rho g / [\sigma]$；
	\item 离散化 $x$ 坐标，计算面积分布 $A(x) = A_0 e^{kx}$；
	\item 积分得总体积 $V = \int_{0}^{L} A(x) dx$ 和总重量 $W = \rho g V$；
	\item 验证各截面应力：采用数值积分计算累积重量，验证 $\sigma(x) \equiv [\sigma]$。
\end{enumerate}

\subsection{参数敏感性分析}

改变许用应力 $[\sigma]$ 的值，观察面积分布的变化趋势：
\begin{itemize}
	\item $[\sigma]$ 越大，$A_0$ 越小、$k$ 越小，总截面积减小；
	\item 高强度材料允许更纤细的设计，但需考虑稳定性（屈曲）约束。
\end{itemize}

\section{案例三：多刚体系统的动力学仿真}

\subsection{曲柄-连杆机构}

曲柄-连杆机构是经典的多刚体系统实例，其动力学方程为高度非线性的微分-代数方程组（DAE）。

位置参数间满足闭环几何约束 $\Phi(\mathbf{q}) = 0$，运动方程写为 Lagrange 乘子形式：
\begin{myformula}
	\begin{cases}
		M \ddot{\mathbf{q}} + \Phi_{\mathbf{q}}^T \boldsymbol{\lambda} = \mathbf{Q}(\mathbf{q}, \dot{\mathbf{q}}, t) \\
		\Phi(\mathbf{q}) = \mathbf{0}
	\end{cases}
\end{myformula}
其中 $\boldsymbol{\lambda}$ 为 Lagrange 乘子（约束反力），$\Phi_{\mathbf{q}}$ 为约束 Jacobi 矩阵。

\subsection{数值求解挑战}

DAE 的数值求解比 ODE 更为困难，主要挑战包括：
\begin{enumerate}
	\item \textbf{约束漂移}：数值误差导致约束方程逐渐不满足，需使用 Baumgarte 稳定化或投影法；
	\item \textbf{指标问题}：原始 DAE 为指标-3，需降指标至指标-1 才能有效数值积分；
	\item \textbf{初始条件相容性}：初始条件必须满足约束方程及其导数。
\end{enumerate}

\subsection{降指标方法}

对约束方程 $\Phi(\mathbf{q}) = 0$ 求两次时间导数：
\[
\Phi_{\mathbf{q}} \ddot{\mathbf{q}} = -\left( \Phi_{\mathbf{q}} \dot{\mathbf{q}} \right)_{\mathbf{q}} \dot{\mathbf{q}} - 2 \Phi_{\mathbf{q} t} \dot{\mathbf{q}} - \Phi_{tt} \equiv \boldsymbol{\gamma}
\]
联立运动方程和加速度级约束得指标-1 系统：
\begin{myformula}
	\begin{pmatrix} M & \Phi_{\mathbf{q}}^T \\ \Phi_{\mathbf{q}} & 0 \end{pmatrix}
	\begin{pmatrix} \ddot{\mathbf{q}} \\ \boldsymbol{\lambda} \end{pmatrix}
	=
	\begin{pmatrix} \mathbf{Q} \\ \boldsymbol{\gamma} \end{pmatrix}
\end{myformula}

这一线性系统可在每一步数值积分中求解，获得广义加速度 $\ddot{\mathbf{q}}$ 和约束反力 $\boldsymbol{\lambda}$。

\section{仿真实践要点总结}

\subsection{代码组织结构}

良好的仿真代码应遵循以下结构：
\begin{verbatim}
% 1. 清空环境
clear; clc; close all;

% 2. 定义参数
params.a = 1;  params.y0 = 8;

% 3. 定义微分方程（匿名函数/函数文件）
odefun = @(t, y, a) -a * y;

% 4. 调用求解器
[t, y] = ode45(@(t,y) odefun(t,y,params.a), [0 10], params.y0);

% 5. 后处理与可视化
plot(t, y);  xlabel('t');  ylabel('y(t)');
\end{verbatim}

\subsection{常见调试问题}

\begin{enumerate}
	\item \textbf{步长过小/积分时间过长}：检查是否存在刚性，尝试 \texttt{ode15s}；
	\item \textbf{解发散}：检查物理参数符号（如刚度、阻尼是否为正值），或使用隐式方法；
	\item \textbf{不满足预期精度}：减小容差（\texttt{odeset} 设置 \texttt{RelTol} 和 \texttt{AbsTol}）；
	\item \textbf{事件检测}：使用 \texttt{odeset} 的 \texttt{Events} 选项处理不连续边界。
\end{enumerate}

\subsection{建模与仿真的最佳实践}

\begin{myTheorem}{工程系统动力学仿真的黄金法则}
	\begin{enumerate}
		\item \textbf{从简到繁}：先建立最简单的模型，验证后再逐步增加复杂度；
		\item \textbf{验证先行}：凡是有解析解的情形，先用数值解验证之；
		\item \textbf{参数敏感性}：分析关键参数对结果的影响，识别主导因素；
		\item \textbf{单位一致性}：保持所有物理量的单位体系一致（推荐 SI 单位制）；
		\item \textbf{可视化诊断}：绘制相图、频谱、能量等辅助判断解的正确性；
		\item \textbf{复现性}：记录参数设置、求解器选项、容差等确保结果可复现。
	\end{enumerate}
\end{myTheorem}

\begin{remark}
	系统动力学建模与仿真的精髓不在于掌握具体的编程语言或求解器命令，而在于透彻理解物理系统的本质特性，并能够将这种理解转化为正确、高效、鲁棒的计算模型。
\end{remark}
